\section{Cross-Sectoral Regression Results}\label{sec:horse_race}



This appendix reports descriptive cross-sectional regressions for the benchmark calibration. The object of interest is the sectoral ranking generated by the per-sector decomposition of a 10 percentage-point tariff. We consider two tariff shocks: the baseline US tariff on Chinese imports (Shock A) and the symmetric Chinese tariff on US imports (Shock B). The main cross-sectional design stacks observations across both shocks and all four model economies---China, United States, Euro area, and Rest of World---yielding $n = 160$ observations ($2 \times 4 \times 20$ sectors).

All predictors are taken from the benchmark calibration. \textbf{Bilateral trade share} is the steady-state share of the \emph{imposing country's} final expenditure sourced from the \emph{receiving country's} sector~$i$: $\beta_C^{\text{USA,CHN},i}$ for Shock A and $\beta_C^{\text{CHN,USA},i}$ for Shock B. This predictor measures how concentrated the tariff is across sectors. Because $\beta_C^{k,\mathrm{rec},i}$ can be written as $\alpha_i \times \mu_i$, we also split it into \textbf{sector weight} $\alpha_i \equiv \sum_j \beta_C^{k,j,i}$ and \textbf{import penetration} $\mu_i \equiv \beta_C^{k,\mathrm{rec},i}/\alpha_i$. \textbf{IO intensity} is the intermediate-input share $\vartheta_{ki}$, which is country-specific. \textbf{Role} denotes a country's relationship to the shock: \emph{Receiving} (direct tariff target), \emph{Imposing} (tariff setter), or \emph{ThirdParty} (unaffected by the tariff directly). All appendix tables are generated mechanically by \texttt{run\_horse\_race\_appendix.py}.

\subsection{China: Direct Exposure under Shock A}

We first document the within-China cross-sectional pattern under Shock A, which provides the cleanest identification: Chinese sectors are directly hit by US demand destruction, and the bilateral trade share $\beta_C^{\text{USA,CHN},i}$ directly measures the tariff's sectoral intensity. Table~\ref{tab:horse_race_china_uni} reports univariate regressions for China's aggregate GDP contribution and own-sector value added. Bilateral trade share explains 93.1\% of the cross-sectional variation in GDP contributions. Splitting $\beta_C = \alpha \times \mu$ shows that nearly all of this signal comes from the import-penetration margin: $\alpha$ alone has $R^2 = 0.041$ for GDP contributions and $0.003$ for own-sector value added, while $\mu$ alone has $R^2 = 0.872$ and $0.494$, respectively. IO intensity, sector size, and price flexibility all have weak univariate fits.

\input{generated/horse_race_table_univariate}

Table~\ref{tab:horse_race_china_multi} shows that adding IO intensity to bilateral trade share raises the GDP fit by only $\Delta R^2 = 0.002$ (from $0.931$ to $0.933$). The decomposed $\alpha + \mu$ specification reaches $R^2 = 0.900$ for GDP and $0.495$ for own-sector value added, only slightly above $\mu$ alone. The product $\alpha \times \mu$ therefore fits better than the additive decomposition, but the decomposition still makes clear which component matters: in the benchmark Shock A design, cross-sectoral variation comes primarily from partner-specific import penetration rather than from broad sector importance in the imposing country's expenditure basket. The conclusion is narrow: bilateral trade exposure is the first-order predictor of which Chinese sectors matter most in the cross section.

\input{generated/horse_race_table_multivariate}

\subsection{Role-Interacted Analysis: Two Shocks, Four Countries}\label{subsec:horse_race_role}

We now exploit the full 160-observation design. For each shock, we classify each economy by role (Imposing, Receiving, ThirdParty) and ask whether the bilateral trade predictor organizes sectoral responses differently across roles.

\subsubsection*{Within-role R\textsuperscript{2} by country and shock}

Table~\ref{tab:horse_race_role_grid} reports the within-group $R^2$ for bilateral trade as a predictor of each country's aggregate GDP contribution, for all four countries under both shocks ($n=20$ per cell).

\input{generated/horse_race_table_role_grid}

Several patterns stand out.

\textbf{Receiving countries are strongly organized by bilateral trade under both shocks.} China (Receiving, Shock A) has $R^2 = 0.931$ and United States (Receiving, Shock B) has $R^2 = 0.684$. The lower fit for the US under Shock B reflects the much smaller and less variable bilateral trade shares in the CHN$\to$USA direction: China's final consumption shares from US sectors range from near zero to at most $0.07\%$, compared to US shares from Chinese sectors that reach $0.42\%$. Even so, the cross-sectional ranking is meaningfully organized.

\textbf{Imposing countries are also strongly organized.} United States (Imposing, Shock A) has $R^2 = 0.907$---sectors where the US imposed large tariffs (high US import share from China) also experienced the largest US output responses, consistent with import-substitution and trade-balance adjustment. China (Imposing, Shock B) has $R^2 = 0.435$, again lower because the predictor varies less across sectors.

\textbf{Third-party countries reveal an asymmetry between Shock A and Shock B.} For Shock A, the Rest of World has $R^2 = 0.901$, almost as high as the directly exposed countries. This reflects trade diversion: sectors where the US buys heavily from China (high $\beta_C^{\text{USA,CHN},i}$) are exactly the sectors where RoW exporters can substitute for displaced Chinese supply, so the bilateral trade predictor organizes RoW responses. For Shock B, RoW has $R^2 = 0.000$---the predictor $\beta_C^{\text{CHN,USA},i}$ (Chinese imports from US sectors) is too small and too uniform to create sector-level variation in trade diversion opportunities. The Euro area is the exception to Shock A third-party alignment ($R^2 = 0.068$), consistent with the heterogeneous-DCP invoicing mechanism documented in Section~\ref{sec:invoicing}: dollar invoicing severs the link between bilateral trade exposure and EA sectoral outcomes.

\subsubsection*{Role-interacted horse race}

Table~\ref{tab:horse_race_role_interacted} reports the role-interacted horse race on the full $n=160$ stacked sample. The model includes role fixed effects (Imposing, Receiving, ThirdParty), bilateral trade interacted with role, and IO intensity interacted with role.

\input{generated/horse_race_table_role_interacted}

Role fixed effects alone explain 25.4\% of total variance across the 160 observations---a modest baseline that captures mean differences across imposers, receivers, and bystanders regardless of sector. Adding bilateral trade interacted with role raises $R^2$ from 0.254 to 0.897, a gain of 64.2 percentage points. The role-specific trade slopes tell the economic story: the Receiving slope ($\hat{\beta} = -14.93$) is nearly three times larger in magnitude than the Imposing slope ($\hat{\beta} = -5.90$), while the ThirdParty slope is near zero ($\hat{\beta} = +0.53$). Receivers are hit hardest, importer-country sectors see moderate negative responses (reduced activity in tariff-protected import-competing sectors is offset by trade balance improvement), and third parties are approximately unaffected on average.

Adding IO intensity interacted with role contributes $\Delta R^2 = 0.000$. Even with full role heterogeneity allowed, IO intensity adds no explanatory power for the cross-sectional allocation of sectoral responses. The model's shock propagation mechanism is bilateral-trade-dominated at the sectoral level; input-output intensity operates as an aggregate amplifier (documented in Section~\ref{sec:robustness}) rather than as a cross-sectional sorter.

Panel~B of Table~\ref{tab:horse_race_role_interacted} shows how this trade predictor decomposes on the stacked sample. Replacing bilateral trade with $\alpha$ interacted with role raises $R^2$ only from $0.254$ to $0.282$, whereas replacing it with $\mu$ interacted with role raises $R^2$ to $0.807$. Using both $\alpha$ and $\mu$ lifts the fit further to $0.860$, close to but still below the $0.897$ delivered by the direct bilateral-trade regressor. The decomposition therefore preserves the main message: most of the stacked signal comes from the penetration margin $\mu$, while the exact product $\alpha \times \mu$ still contains information that the additive split does not recover perfectly.

\subsubsection*{Alpha versus mu by country and shock}

Table~\ref{tab:horse_race_alpha_mu} reports the same decomposition separately within each country-shock cell.

\input{generated/horse_race_table_alpha_mu}

Under Shock~A, $\mu$ is the dominant component wherever the benchmark bilateral-trade predictor works well. For China as the receiving country, $\mu$ alone explains 87.2\% of the GDP ranking, versus only 4.1\% for $\alpha$. The same pattern appears for the United States as imposing country ($0.789$ versus $0.014$) and for the rest of the world as third party ($0.786$ versus $0.019$). The euro area remains the exception: under heterogeneous DCP neither $\alpha$ nor $\mu$ organizes the EA GDP ranking well, consistent with the invoicing wedge discussed in the main text.

Under Shock~B, the decomposition changes. Chinese imports from US sectors are small and comparatively flat across sectors, so the penetration margin becomes nearly uninformative for China as imposing country ($R^2 = 0.000$) and for the United States as receiving country ($R^2 = 0.016$). In those cells, the broader sector-weight margin $\alpha$ carries more of the signal ($R^2 = 0.218$ for China and $0.349$ for the United States). The full bilateral-trade regressor still fits better than either component on its own, implying that the multiplicative combination matters even when neither piece dominates in isolation. The decomposition is therefore not universal across shock directions, but it is informative about what drives the benchmark results: the strong US$\to$CHN appendix findings are mostly penetration-driven, not weight-driven.

\subsection{Amplification Factor}

We examine whether the model allocates amplification differentially to IO-intensive sectors. Define the amplification factor as $A_i^k \equiv |GDP_i^k| / s_i^{\text{trade}}$, the model-predicted aggregate GDP contribution per unit of bilateral trade exposure. Table~\ref{tab:horse_race_amplification} regresses $A_i^k$ on IO intensity, stacked across three importing countries under Shock A ($n=57$; observations with zero bilateral trade share excluded).

\input{generated/horse_race_table_amplification}

IO intensity explains essentially none of the variation in the amplification factor ($R^2 = 0.016$, $p = 0.342$). The amplification factor itself is highly dispersed (CV $= 3.65$, range $[0.020, 2096]$): a handful of sectors with near-zero bilateral trade shares but nonzero GDP responses via IO linkages produce large ratios. This dispersion is not organized by IO intensity. The result reinforces the main finding: IO linkages matter in aggregate but do not generate systematically larger per-unit-of-trade-exposure responses for IO-intensive sectors.

\subsection{Country and Invoicing Heterogeneity}\label{subsec:horse_race_regime}

Table~\ref{tab:horse_race_regime} reports bilateral trade predictability under the benchmark heterogeneous-DCP regime and alternative PCP invoicing regimes for Shock A, using each country's aggregate GDP contribution.

\input{generated/horse_race_table_regime}

The euro area result is the key cross-regime comparison. Under heterogeneous DCP, $R^2 = 0.068$ ($p = 0.266$); under PCP, $R^2 = 0.930$ ($p < 0.001$). This invoicing channel explains why the EA is the outlier in Table~\ref{tab:horse_race_role_grid}: under heterogeneous DCP, dollar invoicing prevents the exchange-rate pass-through that would otherwise align sectoral responses with bilateral trade exposure.

\subsection{Aggregate Consumption and CPI}

Table~\ref{tab:horse_race_aggregate} extends the cross-sectional analysis to aggregate consumption and CPI for Shock A across all three original importing countries and both invoicing regimes.

\input{generated/horse_race_table_aggregate_outcomes}

For China and the United States, bilateral trade share organizes consumption and CPI contributions as tightly as GDP. The euro area reveals a sharp channel separation: heterogeneous DCP delivers $R^2 = 0.924$ for EA CPI even while $R^2 = 0.068$ for EA GDP. Under PCP, the pattern reverses---GDP aligns with trade exposure ($R^2 = 0.930$) while CPI loses alignment ($R^2 = 0.128$). This is the dual-channel signature of dominant-currency pricing analyzed in Section~\ref{sec:invoicing}.

Table~\ref{tab:horse_race_alpha_mu_aggregate} repeats the Shock~A aggregate-outcome exercise using $\alpha$ and $\mu$ separately.

\input{generated/horse_race_table_alpha_mu_aggregate}

The decomposition lines up with the GDP results above. For China and the United States, $\mu$ alone already explains most of the variation across GDP, consumption, and CPI, with $R^2$ values between $0.789$ and $0.897$, whereas $\alpha$ is close to zero throughout. For the euro area, $\mu$ explains the CPI ranking strongly ($R^2 = 0.874$) and the consumption ranking moderately ($R^2 = 0.455$), while $\alpha$ remains weak. The new aggregate CPI result is therefore also penetration-driven: the heterogeneous-DCP cost-push channel sorts sectors mainly through their exposure to Chinese supply in the imposing country's spending basket, not through generic sector size alone.

\subsection{Weighted Robustness}

Domar-weighting by sectoral size leaves results unchanged for directly exposed economies: weighted $R^2$ is $0.939$ for China and $0.877$ for the United States. The euro area weighted fit remains weak ($R^2 = 0.034$). Unweighted OLS is the baseline specification throughout, as size-weighting does not alter the core ranking result.

\begin{figure}[p]
  \centering
  \includegraphics[width=\linewidth]{figures/Fig_CN_Sectoral_Heatmap.png}
  \caption{United States (top) and China (bottom): own-sector value added and inflation under a 10 pp US tariff on Chinese imports. Left: own-sector value added (on-impact \% deviation). Right: own-sector inflation (pp, four-quarter sum). Sectors are ordered by value-added magnitude in the single-sector decomposition.}
  \label{fig:cn_heatmap}
\end{figure}

\subsection{Structural Scatter Panels}\label{sec:appendix_structural_scatters}

The figures below collect the full $3\times4$ scatter grids referenced in Section~\ref{sec:sectoral_channels}. They are descriptive companions to the appendix regressions in this section: for China and the United States, the bilateral-trade-share column visually tracks the aggregate GDP ranking, whereas the benchmark euro-area mapping is much flatter under heterogeneous DCP. The main text therefore keeps only the summary scatter in Figure~\ref{fig:cn_scatter} and relegates the full grids here.

\begin{figure}[H]
  \centering
  \includegraphics[width=\linewidth]{figures/Fig_EA_Structural_Scatter_3x4.png}
  \caption{Euro area: structural determinants of sectoral importance under a 10 percentage-point US tariff on Chinese imports (indirect exposure). $3\times4$ panel: rows show aggregate GDP and consumption contributions (on impact) and CPI contributions (four-quarter rolling sum); columns show bilateral trade share (US imports from China), consumption share, IO linkage intensity ($\vartheta_{ki}$), and price flexibility ($\kappa_{ki}$). Each panel reports $R^2$ and $p$-value from the univariate regression. Each blue point is one of 20 tradeable EA sectors; red dashed line is linear trend. Each sector's response to a 10 percentage-point tariff increase on that sector alone (per-sector decomposition).}
  \label{fig:cn_scatter_3x3}
\end{figure}

\begin{figure}[H]
  \centering
  \includegraphics[width=\linewidth]{figures/Fig_CHN_Structural_Scatter_3x4.png}
  \caption{China: structural determinants of sectoral importance under a 10 percentage-point US tariff on Chinese imports (direct exposure). $3\times4$ panel; see Figure~\ref{fig:cn_scatter_3x3} caption for variable definitions. Each gold point is one of 20 tradeable Chinese sectors; red dashed line is linear trend. Bilateral trade share dominates the GDP ranking ($R^2 = 0.931$).}
  \label{fig:chn_scatter_3x3}
\end{figure}

\begin{figure}[H]
  \centering
  \includegraphics[width=\linewidth]{figures/Fig_US_Structural_Scatter_3x4.png}
  \caption{United States: structural determinants of sectoral importance under a 10 percentage-point US tariff on Chinese imports. $3\times4$ panel: rows show aggregate GDP, consumption, and CPI contributions; columns show bilateral trade share, consumption share, IO linkage intensity ($\vartheta_{ki}$), and price flexibility ($\kappa_{ki}$). Each panel reports $R^2$ and $p$-value. Each orange point is one of 20 tradeable US sectors; dashed line is linear trend. Bilateral trade share dominates the GDP ranking ($R^2 = 0.907$).}
  \label{fig:us_scatter_3x3}
\end{figure}

\section{Additional Figures and Tables}\label{app:C}
Table \ref{tab:nace_sectors} lists the 44 NACE sectors used in the analysis.

% \newpage
\begin{longtable}{ll}
\caption{NACE Rev. 2 sectors used in the Analysis}
\label{tab:nace_sectors} \\
\toprule
\textbf{NACE Code} & \textbf{Sector Name} \\
\midrule
\endfirsthead

\toprule
\textbf{NACE Code} & \textbf{Sector Name} \\
\midrule
\endhead

% Repeat header on every page
\midrule
\multicolumn{2}{r}{{Continued on next page}} \\
\midrule
\endfoot

\bottomrule
\endlastfoot

A.01--02 & Agriculture, forestry and fishing \\
A.03 & Fishing and aquaculture \\
B.05--06 & Mining of coal, lignite, crude petroleum and natural gas \\
B.07--08 & Mining of metal ores and other mining and quarrying \\
B.09 & Mining support service activities \\
C.10--12 & Manufacture of food, beverages and tobacco \\
C.13--15 & Manufacture of textiles, wearing apparel and leather products \\
C.16 & Manufacture of wood and products of wood and cork \\
C.17--18 & Manufacture of paper, and printing and media reproduction \\
C.19 & Manufacture of coke and refined petroleum products \\
C.20 & Manufacture of chemicals and chemical products \\
C.21 & Manufacture of basic pharmaceutical products \\
C.22 & Manufacture of rubber and plastic products \\
C.23 & Manufacture of other non-metallic mineral products \\
C.24 & Manufacture of basic metals \\
C.25 & Manufacture of fabricated metal products \\
C.26 & Manufacture of computer, electronic and optical products \\
C.27 & Manufacture of electrical equipment \\
C.28 & Manufacture of machinery and equipment n.e.c. \\
C.29 & Manufacture of motor vehicles, trailers and semi-trailers \\
C.30 & Manufacture of other transport equipment \\
C.31--33 & Other manufacturing and repair and installation of machinery \\
D.35 & Electricity, gas, steam and air conditioning supply \\
E.36--39 & Water supply; sewerage, waste management and remediation \\
F.41--43 & Construction \\
G.45--47 & Wholesale and retail trade; repair of motor vehicles \\
H.49 & Land transport and transport via pipelines \\
H.50 & Water transport \\
H.51 & Air transport \\
H.52 & Warehousing and support activities for transportation \\
H.53 & Postal and courier activities \\
I.55--56 & Accommodation and food service activities \\
J.58--60 & Publishing, audiovisual and broadcasting activities \\
J.61 & Telecommunications \\
J.62--63 & IT and other information services \\
K.64--66 & Financial and insurance activities \\
L.68 & Real estate activities \\
M.69--75 & Professional, scientific and technical activities \\
N.77--82 & Administrative and support service activities \\
O.84 & Public administration and defence; compulsory social security \\
P.85 & Education \\
Q.86--88 & Human health and social work activities \\
R.90--93 & Arts, entertainment and recreation \\
S.94--96 & Other personal service activities \\
\end{longtable}

\FloatBarrier
